\documentclass[10pt,a4paper]{amsart}
\usepackage{amsmath,amssymb,amsthm,graphicx}
% For commutative diagrams we use Xy-pic
\usepackage[all]{xy}

% We use 2cell for 2-commutative diagrams.
\xyoption{2cell}
% Theorem environments.
%
\theoremstyle{plain}
\newtheorem{theorem}[subsection]{Theorem}
\newtheorem{proposition}[subsection]{Proposition}
\newtheorem{lemma}[subsection]{Lemma}

\theoremstyle{definition}
\newtheorem{definition}[subsection]{Definition}
\newtheorem{example}[subsection]{Example}
\newtheorem{exercise}[subsection]{Exercise}
\newtheorem{situation}[subsection]{Situation}

\theoremstyle{remark}
\newtheorem{remark}[subsection]{Remark}
\newtheorem{remarks}[subsection]{Remarks}

\numberwithin{equation}{subsection}

% Macros
%
\def\lim{\mathop{\mathrm{lim}}\nolimits}
\def\colim{\mathop{\mathrm{colim}}\nolimits}
\def\Spec{\mathop{\mathrm{Spec}}}
\def\Hom{\mathop{\mathrm{Hom}}\nolimits}
\def\Ext{\mathop{\mathrm{Ext}}\nolimits}
\def\SheafHom{\mathop{\mathcal{H}\!\mathit{om}}\nolimits}
\def\SheafExt{\mathop{\mathcal{E}\!\mathit{xt}}\nolimits}
\def\Sch{\mathit{Sch}}
\def\Mor{\mathop{\mathrm{Mor}}\nolimits}
\def\Ob{\mathop{\mathrm{Ob}}\nolimits}
\def\Sh{\mathop{\mathit{Sh}}\nolimits}
\def\NL{\mathop{N\!L}\nolimits}
\def\CH{\mathop{\mathrm{CH}}\nolimits}
\def\proetale{{pro\text{-}\acute{e}tale}}
\def\etale{{\acute{e}tale}}
\def\QCoh{\mathit{QCoh}}
\def\Ker{\mathop{\mathrm{Ker}}}
\def\Im{\mathop{\mathrm{Im}}}
\def\Coker{\mathop{\mathrm{Coker}}}
\def\Coim{\mathop{\mathrm{Coim}}}

% Boxtimes
%
\DeclareMathSymbol{\boxtimes}{\mathbin}{AMSa}{"02}

%
% Macros for moduli stacks/spaces
%
\def\QCohstack{\mathcal{QC}\!\mathit{oh}}
\def\Cohstack{\mathcal{C}\!\mathit{oh}}
\def\Spacesstack{\mathcal{S}\!\mathit{paces}}
\def\Quotfunctor{\mathrm{Quot}}
\def\Hilbfunctor{\mathrm{Hilb}}
\def\Curvesstack{\mathcal{C}\!\mathit{urves}}
\def\Polarizedstack{\mathcal{P}\!\mathit{olarized}}
\def\Complexesstack{\mathcal{C}\!\mathit{omplexes}}
% \Pic is the operator that assigns to X its picard group, usage \Pic(X)
% \Picardstack_{X/B} denotes the Picard stack of X over B
% \Picardfunctor_{X/B} denotes the Picard functor of X over B
\def\Pic{\mathop{\mathrm{Pic}}\nolimits}
\def\Picardstack{\mathcal{P}\!\mathit{ic}}
\def\Picardfunctor{\mathrm{Pic}}
\def\Deformationcategory{\mathcal{D}\!\mathit{ef}}
\newenvironment{reference}{}{}
\newenvironment{slogan}{}{}
\newenvironment{history}{}{}
\usepackage{cleveref}
\crefname{originalsection}{original section}{original sections}
\Crefname{originalsection}{Original section}{Original sections}
\makeatletter
\newlabel{section-axioms-functors}{{07XZ}{source}{}{sourcepointer-1}{}}
\newlabel{section-axioms-functors@cref}{{[originalsection][0][]07XZ}{[1][0][]source}{}{}{}}
\makeatother
\crefname{originaldefinition}{original definition}{original definitions}
\Crefname{originaldefinition}{Original definition}{Original definitions}
\makeatletter
\newlabel{definition-versal}{{07XF}{source}{}{sourcepointer-2}{}}
\newlabel{definition-versal@cref}{{[originaldefinition][0][]07XF}{[1][0][]source}{}{}{}}
\makeatother
\crefname{originallemma}{original lemma}{original lemmas}
\Crefname{originallemma}{Original lemma}{Original lemmas}
\makeatletter
\newlabel{lemma-diagonal}{{0CXI}{source}{}{sourcepointer-3}{}}
\newlabel{lemma-diagonal@cref}{{[originallemma][0][]0CXI}{[1][0][]source}{}{}{}}
\makeatother
\crefname{originallemma}{original lemma}{original lemmas}
\Crefname{originallemma}{Original lemma}{Original lemmas}
\makeatletter
\newlabel{algebraic-lemma-representable-diagonal}{{045G}{source}{}{sourcepointer-4}{}}
\newlabel{algebraic-lemma-representable-diagonal@cref}{{[originallemma][0][]045G}{[1][0][]source}{}{}{}}
\makeatother
\crefname{originaldefinition}{original definition}{original definitions}
\Crefname{originaldefinition}{Original definition}{Original definitions}
\makeatletter
\newlabel{spaces-morphisms-definition-finite-type-point}{{06EG}{source}{}{sourcepointer-5}{}}
\newlabel{spaces-morphisms-definition-finite-type-point@cref}{{[originaldefinition][0][]06EG}{[1][0][]source}{}{}{}}
\makeatother
\crefname{originallemma}{original lemma}{original lemmas}
\Crefname{originallemma}{Original lemma}{Original lemmas}
\makeatletter
\newlabel{lemma-fibre-product-deformation-categories}{{07WL}{source}{}{sourcepointer-6}{}}
\newlabel{lemma-fibre-product-deformation-categories@cref}{{[originallemma][0][]07WL}{[1][0][]source}{}{}{}}
\makeatother
\crefname{originallemma}{original lemma}{original lemmas}
\Crefname{originallemma}{Original lemma}{Original lemmas}
\makeatletter
\newlabel{lemma-versal-change-of-field}{{07XG}{source}{}{sourcepointer-7}{}}
\newlabel{lemma-versal-change-of-field@cref}{{[originallemma][0][]07XG}{[1][0][]source}{}{}{}}
\makeatother
\crefname{originallemma}{original lemma}{original lemmas}
\Crefname{originallemma}{Original lemma}{Original lemmas}
\makeatletter
\newlabel{formal-defos-lemma-smooth-properties}{{06HM}{source}{}{sourcepointer-8}{}}
\newlabel{formal-defos-lemma-smooth-properties@cref}{{[originallemma][0][]06HM}{[1][0][]source}{}{}{}}
\makeatother
\crefname{originallemma}{original lemma}{original lemmas}
\Crefname{originallemma}{Original lemma}{Original lemmas}
\makeatletter
\newlabel{lemma-formally-smooth-on-deformation-categories}{{07WK}{source}{}{sourcepointer-9}{}}
\newlabel{lemma-formally-smooth-on-deformation-categories@cref}{{[originallemma][0][]07WK}{[1][0][]source}{}{}{}}
\makeatother
\crefname{originallemma}{original lemma}{original lemmas}
\Crefname{originallemma}{Original lemma}{Original lemmas}
\makeatletter
\newlabel{more-morphisms-lemma-lifting-along-artinian-at-point}{{02HX}{source}{}{sourcepointer-10}{}}
\newlabel{more-morphisms-lemma-lifting-along-artinian-at-point@cref}{{[originallemma][0][]02HX}{[1][0][]source}{}{}{}}
\makeatother
\crefname{originalsection}{original section}{original sections}
\Crefname{originalsection}{Original section}{Original sections}
\makeatletter
\newlabel{situation-contractions}{{0GH8}{source}{}{sourcepointer-11}{}}
\newlabel{situation-contractions@cref}{{[originalsection][0][]0GH8}{[1][0][]source}{}{}{}}
\makeatother
\crefname{originallemma}{original lemma}{original lemmas}
\Crefname{originallemma}{Original lemma}{Original lemmas}
\makeatletter
\newlabel{lemma-algebraic-stack-RS}{{07WQ}{source}{}{sourcepointer-12}{}}
\newlabel{lemma-algebraic-stack-RS@cref}{{[originallemma][0][]07WQ}{[1][0][]source}{}{}{}}
\makeatother
\crefname{originallemma}{original lemma}{original lemmas}
\Crefname{originallemma}{Original lemma}{Original lemmas}
\makeatletter
\newlabel{formal-spaces-lemma-structure-locally-noetherian}{{0GHM}{source}{}{sourcepointer-13}{}}
\newlabel{formal-spaces-lemma-structure-locally-noetherian@cref}{{[originallemma][0][]0GHM}{[1][0][]source}{}{}{}}
\makeatother
\crefname{originalsection}{original section}{original sections}
\Crefname{originalsection}{Original section}{Original sections}
\makeatletter
\newlabel{formal-spaces-section-adic}{{0AQ2}{source}{}{sourcepointer-14}{}}
\newlabel{formal-spaces-section-adic@cref}{{[originalsection][0][]0AQ2}{[1][0][]source}{}{}{}}
\makeatother
\crefname{originaldefinition}{original definition}{original definitions}
\Crefname{originaldefinition}{Original definition}{Original definitions}
\makeatletter
\newlabel{bootstrap-definition-property-transformation}{{03XZ}{source}{}{sourcepointer-15}{}}
\newlabel{bootstrap-definition-property-transformation@cref}{{[originaldefinition][0][]03XZ}{[1][0][]source}{}{}{}}
\makeatother
\crefname{originallemma}{original lemma}{original lemmas}
\Crefname{originallemma}{Original lemma}{Original lemmas}
\makeatletter
\newlabel{formal-spaces-lemma-diagonal-morphism-formal-algebraic-spaces}{{0AN2}{source}{}{sourcepointer-16}{}}
\newlabel{formal-spaces-lemma-diagonal-morphism-formal-algebraic-spaces@cref}{{[originallemma][0][]0AN2}{[1][0][]source}{}{}{}}
\makeatother
\crefname{originallemma}{original lemma}{original lemmas}
\Crefname{originallemma}{Original lemma}{Original lemmas}
\makeatletter
\newlabel{restricted-lemma-algebraize-rig-etale-affine}{{0AR4}{source}{}{sourcepointer-17}{}}
\newlabel{restricted-lemma-algebraize-rig-etale-affine@cref}{{[originallemma][0][]0AR4}{[1][0][]source}{}{}{}}
\makeatother
\crefname{originallemma}{original lemma}{original lemmas}
\Crefname{originallemma}{Original lemma}{Original lemmas}
\makeatletter
\newlabel{algebra-lemma-what-does-it-mean-again}{{0AS8}{source}{}{sourcepointer-18}{}}
\newlabel{algebra-lemma-what-does-it-mean-again@cref}{{[originallemma][0][]0AS8}{[1][0][]source}{}{}{}}
\makeatother
\crefname{originallemma}{original lemma}{original lemmas}
\Crefname{originallemma}{Original lemma}{Original lemmas}
\makeatletter
\newlabel{algebra-lemma-flat-fibre-smooth}{{00TF}{source}{}{sourcepointer-19}{}}
\newlabel{algebra-lemma-flat-fibre-smooth@cref}{{[originallemma][0][]00TF}{[1][0][]source}{}{}{}}
\makeatother
\title{One common smooth neighbourhood for an adic versal formal morphism}
\author{The Stacks Project Authors}
\date{stacks-project:artin.tex:local-snapshot-20261011; GNU FDL 1.2 or later}
\begin{document}
\maketitle
\noindent\textit{Editorial note.} The final source sentence reverses the ring arrow: the asserted smooth map is B\_n to A\_n, as in all preceding definitions, not A\_n to B\_n.\par\medskip
\noindent\textit{Editorial note.} The source calls the associated graded rings Rees algebras; the displayed direct sums of successive quotients are the actual objects used.\par\medskip
\noindent\textit{Editorial note.} The affine reduction sentence saying V-prime\_n to W-prime\_n is smooth means at the selected point before shrinking; the following sentence explicitly states the need for one common neighbourhood.\par\medskip
\noindent\textit{Allowed standard prerequisites (editorial).}\begin{itemize}
\item Explicit prerequisite boundary. Locally Noetherian formal algebraic spaces and adic maps, Artin versality via square-zero lifting, algebraization of etale maps of affine completions, smoothness descent, finite-module support and the graded flatness criterion.
\end{itemize}\medskip
\ifdefined\hypertarget\hypertarget{sourcepointer-1}{}\fi
\noindent\textit{Source reference.} Original section 07XZ: focused-seeds/AGP/sources/stacks-snapshot-20261011/tags.txt, source lines 10094--10094. Stable Stacks tag for an expressly allowed foundational or notational input; the counted construction remains the two concrete core obligations. The retained local core is not replaced by this pointer. See standard\_prerequisites.\par\medskip
\ifdefined\hypertarget\hypertarget{sourcepointer-2}{}\fi
\noindent\textit{Source reference.} Original definition 07XF: focused-seeds/AGP/sources/stacks-snapshot-20261011/tags.txt, source lines 10075--10075. Stable Stacks tag for an expressly allowed foundational or notational input; the counted construction remains the two concrete core obligations. The retained local core is not replaced by this pointer. See standard\_prerequisites.\par\medskip
\ifdefined\hypertarget\hypertarget{sourcepointer-3}{}\fi
\noindent\textit{Source reference.} Original lemma 0CXI: focused-seeds/AGP/sources/stacks-snapshot-20261011/tags.txt, source lines 16208--16208. Stable Stacks tag for an expressly allowed foundational or notational input; the counted construction remains the two concrete core obligations. The retained local core is not replaced by this pointer. See standard\_prerequisites.\par\medskip
\ifdefined\hypertarget\hypertarget{sourcepointer-4}{}\fi
\noindent\textit{Source reference.} Original lemma 045G: focused-seeds/AGP/sources/stacks-snapshot-20261011/tags.txt, source lines 5435--5435. Stable Stacks tag for an expressly allowed foundational or notational input; the counted construction remains the two concrete core obligations. The retained local core is not replaced by this pointer. See standard\_prerequisites.\par\medskip
\ifdefined\hypertarget\hypertarget{sourcepointer-5}{}\fi
\noindent\textit{Source reference.} Original definition 06EG: focused-seeds/AGP/sources/stacks-snapshot-20261011/tags.txt, source lines 8215--8215. Stable Stacks tag for an expressly allowed foundational or notational input; the counted construction remains the two concrete core obligations. The retained local core is not replaced by this pointer. See standard\_prerequisites.\par\medskip
\ifdefined\hypertarget\hypertarget{sourcepointer-6}{}\fi
\noindent\textit{Source reference.} Original lemma 07WL: focused-seeds/AGP/sources/stacks-snapshot-20261011/tags.txt, source lines 10046--10046. Stable Stacks tag for an expressly allowed foundational or notational input; the counted construction remains the two concrete core obligations. The retained local core is not replaced by this pointer. See standard\_prerequisites.\par\medskip
\ifdefined\hypertarget\hypertarget{sourcepointer-7}{}\fi
\noindent\textit{Source reference.} Original lemma 07XG: focused-seeds/AGP/sources/stacks-snapshot-20261011/tags.txt, source lines 10076--10076. Stable Stacks tag for an expressly allowed foundational or notational input; the counted construction remains the two concrete core obligations. The retained local core is not replaced by this pointer. See standard\_prerequisites.\par\medskip
\ifdefined\hypertarget\hypertarget{sourcepointer-8}{}\fi
\noindent\textit{Source reference.} Original lemma 06HM: focused-seeds/AGP/sources/stacks-snapshot-20261011/tags.txt, source lines 8326--8326. Stable Stacks tag for an expressly allowed foundational or notational input; the counted construction remains the two concrete core obligations. The retained local core is not replaced by this pointer. See standard\_prerequisites.\par\medskip
\ifdefined\hypertarget\hypertarget{sourcepointer-9}{}\fi
\noindent\textit{Source reference.} Original lemma 07WK: focused-seeds/AGP/sources/stacks-snapshot-20261011/tags.txt, source lines 10045--10045. Stable Stacks tag for an expressly allowed foundational or notational input; the counted construction remains the two concrete core obligations. The retained local core is not replaced by this pointer. See standard\_prerequisites.\par\medskip
\ifdefined\hypertarget\hypertarget{sourcepointer-10}{}\fi
\noindent\textit{Source reference.} Original lemma 02HX: focused-seeds/AGP/sources/stacks-snapshot-20261011/tags.txt, source lines 3280--3280. Stable Stacks tag for an expressly allowed foundational or notational input; the counted construction remains the two concrete core obligations. The retained local core is not replaced by this pointer. See standard\_prerequisites.\par\medskip
\ifdefined\hypertarget\hypertarget{sourcepointer-11}{}\fi
\noindent\textit{Source reference.} Original section 0GH8: focused-seeds/AGP/sources/stacks-snapshot-20261011/tags.txt, source lines 20573--20573. Stable Stacks tag for an expressly allowed foundational or notational input; the counted construction remains the two concrete core obligations. The retained local core is not replaced by this pointer. See standard\_prerequisites.\par\medskip
\ifdefined\hypertarget\hypertarget{sourcepointer-12}{}\fi
\noindent\textit{Source reference.} Original lemma 07WQ: focused-seeds/AGP/sources/stacks-snapshot-20261011/tags.txt, source lines 10050--10050. Stable Stacks tag for an expressly allowed foundational or notational input; the counted construction remains the two concrete core obligations. The retained local core is not replaced by this pointer. See standard\_prerequisites.\par\medskip
\ifdefined\hypertarget\hypertarget{sourcepointer-13}{}\fi
\noindent\textit{Source reference.} Original lemma 0GHM: focused-seeds/AGP/sources/stacks-snapshot-20261011/tags.txt, source lines 20587--20587. Stable Stacks tag for an expressly allowed foundational or notational input; the counted construction remains the two concrete core obligations. The retained local core is not replaced by this pointer. See standard\_prerequisites.\par\medskip
\ifdefined\hypertarget\hypertarget{sourcepointer-14}{}\fi
\noindent\textit{Source reference.} Original section 0AQ2: focused-seeds/AGP/sources/stacks-snapshot-20261011/tags.txt, source lines 13495--13495. Stable Stacks tag for an expressly allowed foundational or notational input; the counted construction remains the two concrete core obligations. The retained local core is not replaced by this pointer. See standard\_prerequisites.\par\medskip
\ifdefined\hypertarget\hypertarget{sourcepointer-15}{}\fi
\noindent\textit{Source reference.} Original definition 03XZ: focused-seeds/AGP/sources/stacks-snapshot-20261011/tags.txt, source lines 5166--5166. Stable Stacks tag for an expressly allowed foundational or notational input; the counted construction remains the two concrete core obligations. The retained local core is not replaced by this pointer. See standard\_prerequisites.\par\medskip
\ifdefined\hypertarget\hypertarget{sourcepointer-16}{}\fi
\noindent\textit{Source reference.} Original lemma 0AN2: focused-seeds/AGP/sources/stacks-snapshot-20261011/tags.txt, source lines 13425--13425. Stable Stacks tag for an expressly allowed foundational or notational input; the counted construction remains the two concrete core obligations. The retained local core is not replaced by this pointer. See standard\_prerequisites.\par\medskip
\ifdefined\hypertarget\hypertarget{sourcepointer-17}{}\fi
\noindent\textit{Source reference.} Original lemma 0AR4: focused-seeds/AGP/sources/stacks-snapshot-20261011/tags.txt, source lines 13532--13532. Stable Stacks tag for an expressly allowed foundational or notational input; the counted construction remains the two concrete core obligations. The retained local core is not replaced by this pointer. See standard\_prerequisites.\par\medskip
\ifdefined\hypertarget\hypertarget{sourcepointer-18}{}\fi
\noindent\textit{Source reference.} Original lemma 0AS8: focused-seeds/AGP/sources/stacks-snapshot-20261011/tags.txt, source lines 13571--13571. Stable Stacks tag for an expressly allowed foundational or notational input; the counted construction remains the two concrete core obligations. The retained local core is not replaced by this pointer. See standard\_prerequisites.\par\medskip
\ifdefined\hypertarget\hypertarget{sourcepointer-19}{}\fi
\noindent\textit{Source reference.} Original lemma 00TF: focused-seeds/AGP/sources/stacks-snapshot-20261011/tags.txt, source lines 1073--1073. Stable Stacks tag for an expressly allowed foundational or notational input; the counted construction remains the two concrete core obligations. The retained local core is not replaced by this pointer. See standard\_prerequisites.\par\medskip
\begin{lemma}
\label{lemma-openness-smoothness}
Let $S$ be a locally Noetherian scheme. Let $V$ be a scheme locally
of finite type over $S$. Let $Z \subset V$ be closed. Let $W$ be
a locally Noetherian formal algebraic space over $S$ such that
$W_{red}$ is locally of finite type over $S$. Let $g : V_{/Z} \to W$
be an adic morphism of formal algebraic spaces over $S$. Let $v \in V$
be a closed point such that $g$ is versal at $v$ (as in
Section \ref{section-axioms-functors}).
Then after replacing $V$ by an open neighbourhood of $v$ the
morphism $g$ is smooth (see proof).
\end{lemma}

\begin{lemma}
\label{lemma-base-change-versal}
Let $S$, $\mathcal{X}$, $U$, $x$, $u_0$ be as in
Definition \ref{definition-versal}. Assume
\begin{enumerate}
\item $\Delta : \mathcal{X} \to \mathcal{X} \times \mathcal{X}$
is representable by algebraic spaces,
\item $\Delta$ is locally of finite type
(for example if $\mathcal{X}$ is limit preserving), and
\item $\mathcal{X}$ has (RS).
\end{enumerate}
Let $V$ be a scheme locally of finite type over $S$
and let $y$ be an object of $\mathcal{X}$ over $V$.
Form the $2$-fibre product
$$
\xymatrix{
\mathcal{Z} \ar[r] \ar[d] & (\Sch/U)_{fppf} \ar[d]^x \\
(\Sch/V)_{fppf} \ar[r]^y & \mathcal{X}
}
$$
Let $Z$ be the algebraic space representing $\mathcal{Z}$
and let $z_0 \in |Z|$ be a finite type point lying over $u_0$.
If $x$ is versal at $u_0$, then
the morphism $Z \to V$ is smooth at $z_0$.
\end{lemma}

\begin{proof}
(The parenthetical remark in the statement holds by
Lemma \ref{lemma-diagonal}.)
Observe that $Z$ exists by assumption (1) and Algebraic Stacks, Lemma
\ref{algebraic-lemma-representable-diagonal}. By assumption (2) we see that
$Z \to V \times_S U$ is locally of finite type.
Choose a scheme $W$, a closed point $w_0 \in W$, and
an \'etale morphism $W \to Z$ mapping $w_0$ to $z_0$, see
Morphisms of Spaces, Definition
\ref{spaces-morphisms-definition-finite-type-point}.
Then $W$ is locally of finite type over $S$ and
$w_0$ is a finite type point of $W$.
Let $l = \kappa(z_0)$. Denote $z_{l, 0}$, $v_{l, 0}$,
$u_{l, 0}$, and $x_{l, 0}$ the objects of
$\mathcal{Z}$, $(\Sch/V)_{fppf}$, $(\Sch/U)_{fppf}$,
and $\mathcal{X}$ over $\Spec(l)$ obtained by pullback to $\Spec(l) = w_0$.
Consider
$$
\xymatrix{
\mathcal{F}_{(\Sch/W)_{fppf}, l, w_0} \ar[r] &
\mathcal{F}_{\mathcal{Z}, l, z_{l, 0}} \ar[d] \ar[r] &
\mathcal{F}_{(\Sch/U)_{fppf}, l, u_{l, 0}} \ar[d] \\
& \mathcal{F}_{(\Sch/V)_{fppf}, l, v_{l, 0}} \ar[r] &
\mathcal{F}_{\mathcal{X}, l, x_{l, 0}}
}
$$
By Lemma \ref{lemma-fibre-product-deformation-categories}
the square is a fibre product of predeformation categories.
By Lemma \ref{lemma-versal-change-of-field}
we see that the right vertical arrow is smooth.
By Formal Deformation Theory, Lemma
\ref{formal-defos-lemma-smooth-properties}
the left vertical arrow is smooth.
By Lemma \ref{lemma-formally-smooth-on-deformation-categories}
we see that the left horizontal arrow is smooth.
We conclude that the map
$$
\mathcal{F}_{(\Sch/W)_{fppf}, l, w_0} \to
\mathcal{F}_{(\Sch/V)_{fppf}, l, v_{l, 0}}
$$
is smooth by Formal Deformation Theory, Lemma
\ref{formal-defos-lemma-smooth-properties}.
Thus we conclude that $W \to V$ is smooth at $w_0$ by
More on Morphisms, Lemma
\ref{more-morphisms-lemma-lifting-along-artinian-at-point}.
This exactly means that $Z \to V$ is smooth at $z_0$
and the proof is complete.
\end{proof}

\begin{lemma}
\label{lemma-rs}
In Situation \ref{situation-contractions} the functor $F$ satisfies
the Rim-Schlessinger condition (RS).
\end{lemma}

\begin{proof}
Recall that the condition only involves the evaluation $F(V)$ of the functor
$F$ on schemes $V$ over $S$ which are spectra of Artinian local rings
and the restriction maps $F(V_2) \to F(V_1)$ for morphisms $V_1 \to V_2$
of schemes over $S$ which are spectra of Artinian local rings.
Thus let $V/S$ be the spetruim of an Artinian local ring.
If $\xi = (Z, u', \hat x) \in F(V)$ then either $Z = \emptyset$
or $Z = V$ (set theoretically). In the first case we see that
$\hat x$ is a morphism from the empty formal algebraic space
into $W$. In the second case we see that $u'$ is a morphism from
the empty scheme into $X'$ and we see that $\hat x : V \to W$
is a morphism into $W$. We conclude that
$$
F(V) = U'(V) \amalg W(V)
$$
and moreover for $V_1 \to V_2$ as above the induced map
$F(V_2) \to F(V_1)$ is compatible with this decomposition.
Hence it suffices to prove that both $U'$ and $W$ satisfy the
Rim-Schlessinger condition. For $U'$ this follows from
Lemma \ref{lemma-algebraic-stack-RS}.
To see that it is true for $W$, we write $W = \colim W_n$ as in
Formal Spaces, Lemma
\ref{formal-spaces-lemma-structure-locally-noetherian}.
Say $V = \Spec(A)$ with $(A, \mathfrak m)$ an Artinian local ring.
Pick $n \geq 1$ such that $\mathfrak m^n = 0$. Then we have
$W(V) = W_n(V)$. Hence we see that the Rim-Schlessinger condition
for $W$ follows from the Rim-Schlessinger condition for $W_n$ for
all $n$ (which in turn follows from
Lemma \ref{lemma-algebraic-stack-RS}).
\end{proof}

\begin{proof}
Since $g$ is adic it is representable by algebraic spaces (Formal Spaces,
Section \ref{formal-spaces-section-adic}).
Thus by saying $g$ is smooth we mean that $g$ should be smooth
in the sense of
Bootstrap, Definition \ref{bootstrap-definition-property-transformation}.

\medskip\noindent
Write $W = \colim W_n$ as in Formal Spaces, Lemma
\ref{formal-spaces-lemma-structure-locally-noetherian}.
Set $V_n = V_{/Z} \times_{\hat x, W} W_n$.
Then $V_n$ is a closed subscheme with underlying set $Z$.
Smoothness of $V \to W$ is equivalent to the smoothness
of all the morphisms $V_n \to W_n$ (this holds because any morphism
$T \to W$ with $T$ a quasi-compact scheme factors through
$W_n$ for some $n$). We know that the morphism $V_n \to W_n$
is smooth at $v$ by
Lemma \ref{lemma-base-change-versal}\footnote{The lemma applies since
the diagonal of $W$ is representable by algebraic spaces and
locally of finite type, see Formal Spaces, Lemma
\ref{formal-spaces-lemma-diagonal-morphism-formal-algebraic-spaces}
and we have seen that $W$ has (RS) in the proof of Lemma \ref{lemma-rs}.}.
Of course this means that given any $n$ we can shrink $V$
such that $V_n \to W_n$ is smooth. The problem is to find
an open which works for all $n$ at the same time.

\medskip\noindent
The question is local on $V$, hence we may assume $S = \Spec(R)$ and
$V = \Spec(A)$ are affine.

\medskip\noindent
In this paragraph we reduce to the case where $W$ is an affine formal
algebraic space. Choose an affine formal scheme $W'$ and an \'etale morphism
$W' \to W$ such that the image of $v$ in $W_{red}$ is in the
image of $W'_{red} \to W_{red}$. Then $V_{/Z} \times_{g, W} W' \to V_{/Z}$
is an adic \'etale morphism of formal algebraic spaces over $S$
and $V_{/Z} \times_{g, W} W'$ is an affine formal algebraic space.
By Algebraization of Formal Spaces,
Lemma \ref{restricted-lemma-algebraize-rig-etale-affine}
there exists an \'etale morphism $\varphi : V' \to V$ of affine schemes
such that the completion of $V'$ along $Z' = \varphi^{-1}(Z)$
is isomorphic to $V_{/Z} \times_{g, W} W'$ over $V_{/Z}$.
Observe that $v$ is the image of some $v' \in V'$.
Since smoothness is preserved under base change we see that
$V'_n \to W'_n$ is smooth for all $n$. In the next paragraph
we show that after replacing $V'$ by an open neighbourhood of $v'$
the morphisms $V'_n \to W'_n$ are smooth for all $n$.
Then, after we replace $V$ by the open image of $V' \to V$,
we obtain that $V_n \to W_n$ is smooth by \'etale descent of smoothness.
Some details omitted.

\medskip\noindent
Assume $S = \Spec(R)$, $V = \Spec(A)$, $Z = V(I)$, and $W = \text{Spf}(B)$.
Let $v$ correspond to the maximal ideal $I \subset \mathfrak m \subset A$.
We are given an adic continuous $R$-algebra homomorphism
$$
B \longrightarrow A^\wedge
$$
Let $\mathfrak b \subset B$ be the ideal of topologically nilpotent
elements (this is the maximal ideal of definition of the Noetherian adic
topological ring $B$). Observe that $\mathfrak b A^\wedge$ and
$IA^\wedge$ are both ideals of definition of the Noetherian adic
ring $A^\wedge$. Also, $\mathfrak m A^\wedge$ is a maximal ideal
of $A^\wedge$ containing both $\mathfrak b A^\wedge$ and $IA^\wedge$.
We are given that
$$
B_n = B/\mathfrak b^n \to A^\wedge/\mathfrak b^n A^\wedge = A_n
$$
is smooth at $\mathfrak m$ for all $n$. By the discussion above
we may and do assume that $B_1 \to A_1$ is a smooth ring map.
Denote $\mathfrak m_1 \subset A_1$ the maximal ideal corresponding
to $\mathfrak m$. Since smoothness implies flatness, we see that:
for all $n \geq 1$ the map
$$
\mathfrak b^n/\mathfrak b^{n + 1} \otimes_{B_1} (A_1)_{\mathfrak m_1}
\longrightarrow
\left(\mathfrak b^nA^\wedge/\mathfrak b^{n + 1}A^\wedge\right)_{\mathfrak m_1}
$$
is an isomorphism (see
Algebra, Lemma \ref{algebra-lemma-what-does-it-mean-again}).
Consider the Rees algebra
$$
B' = \bigoplus\nolimits_{n \geq 0} \mathfrak b^n/\mathfrak b^{n + 1}
$$
which is a finite type graded algebra over the Noetherian ring $B_1$ and
the Rees algebra
$$
A' = \bigoplus\nolimits_{n \geq 0}
\mathfrak b^nA^\wedge/\mathfrak b^{n + 1}A^\wedge
$$
which is a a finite type graded algebra over the Noetherian ring $A_1$.
Consider the homomorphism of graded $A_1$-algebras
$$
\Psi : B' \otimes_{B_1} A_1 \longrightarrow A'
$$
By the above this map is an isomorphism after localizing at
the maximal ideal $\mathfrak m_1$ of $A_1$.
Hence $\Ker(\Psi)$, resp.\ $\Coker(\Psi)$ is a finite module
over $B' \otimes_{B_1} A_1$, resp.\ $A'$ whose localization
at $\mathfrak m_1$ is zero. It follows that after replacing
$A_1$ (and correspondingly $A$) by a principal localization
we may assume $\Psi$ is an isomorphism. (This is the key step of the proof.)
Then working backwards we see that $B_n \to A_n$ is flat, see
Algebra, Lemma \ref{algebra-lemma-what-does-it-mean-again}.
Hence $A_n \to B_n$ is smooth (as a flat ring map with smooth
fibres, see Algebra, Lemma \ref{algebra-lemma-flat-fibre-smooth})
and the proof is complete.
\end{proof}

\bibliographystyle{plain}
\bibliography{my}
\end{document}
